\section{Implicit inverse composition and real branch selection}
\label{sec:inverse-function-expressions}

An expression $F^{-1}(h(x))$ poses two distinct mathematical problems:
selecting the inverse branch of $F$, and expanding its composition with
$h$. Branch selection must precede the coefficient calculation. A real
inverse germ consists of a defining equation, a source domain, a source
endpoint, and an approach direction. These data distinguish solutions
which may share the same limiting target.

\subsection{Scalar implicit equations and parameters}

For a function $F$ of $n$ real variables, inversion in its $k$th argument
with the other arguments fixed means solving the scalar equation
\begin{equation}
 F(a_1,\ldots,a_{k-1},t,a_{k+1},\ldots,a_n)=a_k.
 \label{eq:inverse-expression-selected-argument}
\end{equation}
This is a one-dimensional inverse problem with parameters. For example,
$a t+t^2=y$, with fixed $a>0$ and $t\ge0$, has the positive local solution
\[
 t=\frac{\sqrt{a^2+4y}-a}{2}
   =\frac ya-\frac{y^2}{a^3}+\frac{2y^3}{a^5}
     +\OO(y^4),\qquad y\to0^+.
\]
The parameter $a$ is fixed in this statement. If it varies with $y$, the
error bound and even the relevant leading balance can change.

An important family with varying coefficients admits the exact reduction
\[
 A(x)F(t)+B(x)=h(x)
 \quad\Longleftrightarrow\quad
 F(t)=\frac{h(x)-B(x)}{A(x)}.
\]
Assume $A(x)\ne0$ on the selected deleted neighborhood and that the
source domain for $t$ is independent of $x$. Then the variation of $A$
and $B$ belongs entirely to the target argument. Nonvanishing at the
endpoint itself is unnecessary: $A(x)=x$ is permitted near $0^+$.
The quotient must nevertheless be expanded with its full uncertainty;
replacing a varying coefficient by its limiting value is generally invalid.

The coupled family $a(x)t+t^2=h(x)$ cannot in general be reduced to a
fixed forward germ by this separation. Joint implicit asymptotics would
require hypotheses on the relative rates and on uniform regularity in the
varying parameters. Such an additional analysis is independent of the
fixed-parameter inversion formulas.

\subsection{Source, target, and parameter domains}

Let $C(t)$ specify a source restriction. The admissible real source domain
is its intersection with the real domain of the chosen forward expression
$F(t)$, under the fixed parameter assumptions. Merely requiring $t$ real
does not imply $t>0$, whereas $t+t^2(1+\log t)$, with the real logarithm,
is defined only for $t>0$. Domain restrictions remain part of the inverse
problem throughout the calculation.

A restriction on the expansion variable $x$ instead specifies the approach
in $F^{-1}(h(x))$. A condition on a fixed parameter specifies the family
for which the conclusion holds. These three kinds of conditions have
different roles: an approach condition must hold eventually along $x$,
a source condition must hold on the selected inverse germ, and a parameter
condition is fixed throughout the limiting argument.

\subsection{Local existence and unambiguous selection}

Let $D\subset\R$ denote the admitted source domain. A source endpoint
$t_*$ is approached through a positive coordinate $u\to0^+$:
\[
 t=t_*+s u\quad(t_*\in\R),\qquad
 t=\frac{s}{u}\quad(t_*=s\infty),\qquad s\in\{-1,1\}.
\]
The condition $t(u)\in D$ concerns a \emph{deleted} neighborhood
$0<u<\delta$. It need not hold at $u=0$, where the source expression
may be undefined. Thus a restriction $0<t<1$ is compatible with
$t\to0^+$.

\begin{proposition}[A selected real inverse germ]
\label{prop:inverse-expression-local-branch}
Suppose there is a deleted source neighborhood on which $F$ is real,
continuous, and differentiable, and $F'$ has a fixed strict sign.
Suppose also that $F(t(u))\to\eta$, with the prescribed one-sided target
approach when $\eta$ is finite. Then a sufficiently small interval on
that target side, or a sufficiently distant target tail when
$\eta=\pm\infty$, has a unique inverse with values in this source
neighborhood. It approaches $t_*$ on the selected source side.
\end{proposition}

\begin{proof}
The source chart parametrizes an interval. The derivative condition
makes $F$ strictly monotone there, and continuity makes its image an
interval. Its endpoint limit and target side show that the image
contains a sufficiently small target interval or tail of the stated
kind. The inverse on that image is unique and continuous, and its
endpoint behavior follows from strict monotonicity.
\end{proof}

A nonzero leading block of the derivative's real asymptotic expansion
can prove its eventual sign. The leading coefficient and the parity of
the leading logarithmic degree determine that sign because
$\log u\to-\infty$. This is a qualitative eventual-sign proof; it does
not compute a numerical radius on which the derivative bound holds.
A neighborhood known to satisfy the source condition is therefore not
automatically a quantitative neighborhood for the full asymptotic error bound.

Proposition~\ref{prop:inverse-expression-local-branch} verifies a proposed
branch. \emph{Selection} among different source endpoints
requires an additional argument. There are two useful sufficient routes.

\begin{enumerate}[leftmargin=*]
\item \textbf{Complete real fiber and boundary information.}
For a finite target $\eta$, all interior points satisfying $F(t)=\eta$
must be accounted for, together with all finite boundaries of $D$ and
both source infinities. Continuity in the interior forces any finite
interior endpoint of an inverse germ to belong to this fiber. Boundary
limits must be checked separately because $F$ need not be defined there.
Every admissible endpoint and source direction is then tested by
Proposition~\ref{prop:inverse-expression-local-branch}. Exactly one
successful choice gives an unambiguous branch. For an infinite target,
the finite interior fiber is replaced by the observation that a
continuous finite-valued function cannot diverge at an interior point.
\item \textbf{Strict monotonicity on a connected real domain.}
If $D$ is an interval and $F$ is strictly monotone throughout $D$, at
most one real inverse can have values in $D$. A separately verified
local candidate therefore identifies that inverse even without an explicit enumeration of the limiting fiber. This route requires both connectedness
and the global monotonicity proof; local monotonicity near one candidate
alone does not exclude another branch elsewhere.
\end{enumerate}

An arbitrary finite list of candidate endpoints does not establish that
either route has succeeded. For example, the two source sides of $t^2$ give
opposite inverses at $\eta=0^+$ when only realness is required. Restricting
$t>0$ or $t<0$ selects one of them. An explicit source endpoint and
direction specifies a local problem even when the global inverse is ambiguous.
The chosen candidate must still satisfy the domain, limit, target-side,
and monotonicity hypotheses.

\subsection{Transporting the target argument's uncertainty}

After selecting the branch, composition must account for uncertainty in
both the inverse and its argument. The following ordinary power--log
lemma makes the two contributions explicit. Exact exponential carriers
or other scales require their corresponding chart rules from
Section~\ref{sec:series-calculus}.

\begin{lemma}[Two independent composition errors]
\label{lem:inverse-expression-composition-error}
Let $w\to0^+$ and $\M(w)=1+|\log w|$. Suppose the positive local target
coordinate satisfies
\[
 q(w)\sim c w^\lambda,\quad c,\lambda>0,\qquad
 q(w)-\widehat q(w)=\OO(w^\rho\M(w)^d),\quad \rho>\lambda.
\]
Let $G$ be the selected inverse, or a differentiable observable of it,
on this positive target coordinate, and suppose
\[
 G(v)-T(v)=\OO(v^\beta\M(v)^k),\qquad
 |G'(v)|\le C v^{\nu-1}\M(v)^\ell
\]
for sufficiently small $v>0$. Here $\beta,\nu\in\R$ and
$d,k,\ell\ge0$. Then
\begin{equation}
 G(q(w))-T(\widehat q(w))
 =\OO(w^{\lambda\beta}\M(w)^k)
  +\OO(w^{\rho+\lambda(\nu-1)}\M(w)^{d+\ell}).
 \label{eq:inverse-expression-two-errors}
\end{equation}
\end{lemma}

\begin{proof}
The target uncertainty is $o(w^\lambda)$, so $q$, $\widehat q$, and
every point between them are positive and comparable to $w^\lambda$.
Their logarithmic envelopes are comparable to $\M(w)$. Split the
difference as
\[
 [G(q)-G(\widehat q)]+[G(\widehat q)-T(\widehat q)].
\]
Substitution into the stated outer remainder gives the first term in
\eqref{eq:inverse-expression-two-errors}. The mean value theorem and the
derivative bound give the second. The derivative estimate is an
independent hypothesis: no differentiation of a magnitude Big-$O$
statement is used.
\end{proof}

Thus the two candidate power cutoffs are $\lambda\beta$ and
$\rho+\lambda(\nu-1)$. The smaller governs the final power order;
at equal powers the logarithmic envelopes must also be combined.
A vanishing derivative can improve transported input precision, while
a singular derivative can lower it. For an inverse itself, a suitable
derivative estimate can be obtained from the original equation and
$G'=1/(F'\circ G)$ on the selected branch. A larger formal cutoff cannot recover coefficients hidden by the input
argument's uncertainty.
Resonant contributions at a common weight are combined before complete
blocks are retained or discarded.

\subsection{Logarithmic and irrational-power examples}

Consider $F(t)=t+t^2(1+\log t)$ on $t>0$. Its derivative satisfies
\[
 F'(t)=1+t(3+2\log t)\ge1-2e^{-5/2}>0,
\]
because $t(3+2\log t)$ has its global minimum at $t=e^{-5/2}$.
The strictly increasing positive real branch tends to $0$ as $y\to0^+$.
Its inverse begins
\[
 G(y)=y-y^2(1+\log y)
 +y^3(2\log^2y+5\log y+3)
 +\OO(y^4\M(y)^3).
\]
This is reversion of a real source germ, although the inverse has no
ordinary Taylor series at $y=0$.

For $F(t)=t+t^{\sqrt2}$ on $t>0$,
$F'(t)=1+\sqrt2\,t^{\sqrt2-1}>0$. At $0^+$ its inverse has the expansion
\[
 G(y)=y-y^{\sqrt2}+\sqrt2\,y^{2\sqrt2-1}
 -\frac{6-\sqrt2}{2}\,y^{3\sqrt2-2}
 +\OO(y^{4\sqrt2-3}).
\]
At $+\infty$, a different normalization is required because $t^{\sqrt2}$
is dominant. Put $p=\sqrt2$, $z=y^{1/p}$, and $\varepsilon=z^{1-p}$.
Writing $G(y)=zW(\varepsilon)$ gives
\[
 W^p+\varepsilon W=1,\qquad W(0)=1.
\]
The implicit derivative in $W$ at the origin is $p\ne0$, so $W$ has a
convergent ordinary power series near $\varepsilon=0$. Its first terms are
\begin{align*}
 G(y)=y^{1/p}\Bigl(&
 1-\frac{\varepsilon}{p}
 +\frac{3-p}{2p^2}\varepsilon^2\\
 &{}
 -\frac{(p-2)(p-4)}{3p^3}\varepsilon^3
 +\frac{(5-p)(5-2p)(5-3p)}{24p^4}\varepsilon^4
 +\OO(\varepsilon^5)\Bigr),
\end{align*}
where $\varepsilon=y^{(1-p)/p}$.
The same globally increasing real function thus has different natural
asymptotic coordinates at its two endpoints.

An arithmetic expression such as $1+G(x)$ or $G(x)^r$ is an observable
of the inverse. Its defining identity is obtained by reconstructing the
source value $G(x)$ before substituting into $F$. It is not itself an
inverse of $F$. On compatible branches, inversion of $G$ recovers $F$;
composition with a general observable has no such interpretation.

\subsection{Exact identities and the role of the selected root}

A closed form does not remove the need to specify its real branch.
For example, $W_0(y)$ inverts $t e^t$ on $t\ge-1$, while $W_{-1}(y)$
inverts it on $t\le-1$. Their common threshold and their distinct
large-argument or near-zero behavior require the corresponding source
endpoint data. An exact algebraic radical has the same obligation to use
the branch matching the source domain.

A residual evaluated at a reconstructed source root must verify both the
forward equation and the source restriction. Checking a powered observable
alone can lose the source sign. A quantitative interval argument must
establish the defining conditions on its verification interval; an open
source condition excludes an endpoint at equality. A numerical root
satisfying the condition supplies evidence at that point, while a uniform
inverse statement requires the neighborhood hypotheses proved above.
